This section presents five forms of generalization evidence: hardware invariance, scale invariance, stochastic sampling invariance, held-out prompt generalization, and cross-architecture transfer.

\subsection{Hardware Invariance}
\label{sec:hardware}

\paragraph{Design.} Experiment 37 ran the identical refusal detection paradigm (20 harmful + 20 benign prompts, Qwen2.5-7B, 50-token generation) on two GPU platforms: NVIDIA GeForce RTX 3090 (consumer, Ampere architecture, 24~GB VRAM) and NVIDIA H200 (datacenter, Hopper architecture, 141~GB HBM3e).

\paragraph{Results.}

\begin{table}[H]
\centering
\caption{Hardware invariance: feature correlation between RTX 3090 and H200.}
\label{tab:hardware}
\begin{tabular}{lcc}
\toprule
\textbf{Feature} & \textbf{Pearson $r$} & \textbf{Max \% difference} \\
\midrule
Key norm & $> 0.999$ & $< 0.05\%$ \\
Norm per token & $> 0.999$ & $< 0.05\%$ \\
Effective rank & $> 0.999$ & $< 0.05\%$ \\
Rank entropy & $> 0.999$ & $< 0.05\%$ \\
\bottomrule
\end{tabular}
\end{table}

All four aggregate features correlate at $r > 0.999$ between platforms, with maximum per-prompt differences below 0.05\%. Text outputs were identical for 39/40 prompts; one prompt differed by a single token.

\paragraph{Interpretation.} KV-cache geometry is a \emph{model property}, not a hardware property. The key matrices are determined by model weights and input tokens; the GPU performs identical floating-point operations regardless of its architecture. This result means that classifiers trained on consumer GPUs can be deployed on datacenter infrastructure without retraining, and vice versa.

\subsection{Scale Invariance}
\label{sec:scale}

\paragraph{Design.} Experiment 26 tested whether category geometry is preserved across model scale, using Campaign 2's 16-model ladder spanning 0.6B to 70B parameters across 6 architecture families. For each scale group (small: 0.6--1.1B; medium: 7--9B; large: 32--70B), we computed the 13$\times$13 pairwise AUROC matrix between cognitive categories.

\paragraph{Results.} Spearman rank correlations of the 78-element AUROC vectors between scale groups:

\begin{table}[H]
\centering
\caption{Scale invariance: cross-scale correlation of category geometry.}
\label{tab:scale}
\begin{tabular}{lc}
\toprule
\textbf{Scale comparison} & \textbf{Spearman $\rho$} \\
\midrule
Small--Medium & 0.83 \\
Medium--Large & 0.90 \\
Small--Large & 0.85 \\
\bottomrule
\end{tabular}
\end{table}

The category geometry is preserved from 0.6B to 70B parameters. Coding remains the highest-norm category in all 16 models. The slight increase in $\rho$ from small--medium to medium--large suggests that category geometry \emph{stabilizes} as models scale up.

\paragraph{Caveat.} The LARGE group contains only two models (Qwen2.5-32B-4bit and DeepSeek-R1-Distill-Qwen-32B), both quantized. The $\rho$ values involving this group should be interpreted cautiously given $n = 2$ and the potential confound of quantization-specific geometry. Scale invariance is best-supported in the small--medium range.

\paragraph{Qualitative invariants.} Three structural properties hold at every scale:
\begin{enumerate}[leftmargin=*]
    \item Coding is \#1 by norm (all 16 models).
    \item Confabulation $\approx$ creative writing (AUROC $< 0.70$ at every scale).
    \item Grounded facts $\approx$ confabulation (AUROC $< 0.70$ at every scale).
\end{enumerate}

\subsection{Stochastic Sampling Invariance}
\label{sec:stochastic_robust}

All Campaign 3 experiments through Experiment 46 used greedy decoding ($\text{temperature} = 0$), which raises a deployment concern: real-world LLM usage is stochastic. Experiment 49a tested detection under temperature $= 0.7$ ($\text{top-}p = 0.9$) on RunPod RTX 6000 Ada hardware.

All three core detection capabilities survive stochastic sampling: deception AUROC $\approx 0.85$ ($p < 0.001$, 95\% CI $[0.892, 0.951]$), impossibility AUROC $\approx 0.80$ ($p < 0.001$), and refusal AUROC $\approx 0.66$ ($p = 0.001$). Within-run feature variation is small (CV: honest 0.3\%, deceptive 1.6\%), confirming that stochastic token selection introduces minimal geometric noise.

The degradation is non-uniform: greedy decoding inflates refusal detection most (AUROC $0.91 \to 0.66$) and deception detection least (${\sim}1.00 \to 0.85$). This is consistent with the per-head anatomy (\Cref{sec:perhead}): deception's globally distributed signal is robust to sampling variation, while refusal's concentrated gatekeeper-head signal is more sensitive to the specific token sequence. Practical deployment should train and evaluate under stochastic conditions to avoid overestimating suppression detection.

\subsection{Held-Out Prompt Generalization}
\label{sec:heldout_robust}

Cross-validated AUROC within a fixed prompt battery may overestimate generalization if prompts share systematic features. Experiment 49d tested classifiers on completely new prompts never seen during training: deception AUROC $\approx 0.81$, refusal AUROC $\approx 0.76$, impossibility AUROC $\approx 0.77$.

The held-out degradation is moderate and consistent across conditions ($\Delta \approx 0.15$--$0.19$), confirming that classifiers learn generalizable suppression geometry rather than prompt-specific artifacts. Combined with stochastic sampling invariance, these results establish that KV-cache detection generalizes across both input distribution (new prompts) and output distribution (stochastic tokens).

\subsection{Cross-Architecture Transfer}
\label{sec:transfer}

The cross-model deception transfer results were presented in \Cref{sec:deception}. Here we summarize the methodological insight.

\paragraph{Why LR beats RF for transfer.} Random forest classifiers achieve perfect within-model AUROC (1.0000) but degrade to 0.67 on cross-model transfer. Logistic regression achieves 0.863. The explanation lies in how each classifier handles architecture-specific feature scaling.

RF partitions the feature space with axis-aligned splits that are calibrated to absolute feature magnitudes. When Qwen produces key norms of ${\sim}29{,}000$ and Llama produces ${\sim}19{,}000$, RF splits learned on Qwen misclassify Llama inputs. LR's linear decision boundary, applied after standard scaling, captures the \emph{direction} of feature perturbation (deceptive processing increases norm, rank, and entropy while decreasing norm-per-token) regardless of architecture-specific magnitudes.

This suggests that the output suppression signal is \emph{directionally universal}: all architectures perturb in the same direction when suppressing output, but with different magnitudes. A linear classifier captures the direction; a nonlinear classifier overfits to the magnitude.
